\documentclass[10pt,a4paper]{amsart}
\usepackage{amsmath,amssymb,amsthm,graphicx}
% Retained preamble for the Stacks Project chapter divisors.tex (Divisors).
% The chapter file pinned-original-source.tex carries only \input{preamble} on its first line, so no
% preamble text lives inside the chapter itself.  The Stacks Project uses one single shared preamble for
% every chapter; the snapshot stages it verbatim in sources/stacks-chow/preamble.tex.  The block below is
% copied mechanically, line for line, from that file's lines 14-18 (the Xy-pic setup) and lines 154-217
% (the native theorem declarations together with the mathematical macros).  Nothing was added, removed or
% redefined.  The wrapper class and the amsmath/amssymb/amsthm/graphicx base packages are supplied by the
% assembler, so the class/\externaldocument part of the shared preamble is not repeated here.
% For commutative diagrams we use Xy-pic
\usepackage[all]{xy}

% We use 2cell for 2-commutative diagrams.
\xyoption{2cell}
% Theorem environments.
%
\theoremstyle{plain}
\newtheorem{theorem}[subsection]{Theorem}
\newtheorem{proposition}[subsection]{Proposition}
\newtheorem{lemma}[subsection]{Lemma}

\theoremstyle{definition}
\newtheorem{definition}[subsection]{Definition}
\newtheorem{example}[subsection]{Example}
\newtheorem{exercise}[subsection]{Exercise}
\newtheorem{situation}[subsection]{Situation}

\theoremstyle{remark}
\newtheorem{remark}[subsection]{Remark}
\newtheorem{remarks}[subsection]{Remarks}

\numberwithin{equation}{subsection}

% Macros
%
\def\lim{\mathop{\mathrm{lim}}\nolimits}
\def\colim{\mathop{\mathrm{colim}}\nolimits}
\def\Spec{\mathop{\mathrm{Spec}}}
\def\Hom{\mathop{\mathrm{Hom}}\nolimits}
\def\Ext{\mathop{\mathrm{Ext}}\nolimits}
\def\SheafHom{\mathop{\mathcal{H}\!\mathit{om}}\nolimits}
\def\SheafExt{\mathop{\mathcal{E}\!\mathit{xt}}\nolimits}
\def\Sch{\mathit{Sch}}
\def\Mor{\mathop{\mathrm{Mor}}\nolimits}
\def\Ob{\mathop{\mathrm{Ob}}\nolimits}
\def\Sh{\mathop{\mathit{Sh}}\nolimits}
\def\NL{\mathop{N\!L}\nolimits}
\def\CH{\mathop{\mathrm{CH}}\nolimits}
\def\proetale{{pro\text{-}\acute{e}tale}}
\def\etale{{\acute{e}tale}}
\def\QCoh{\mathit{QCoh}}
\def\Ker{\mathop{\mathrm{Ker}}}
\def\Im{\mathop{\mathrm{Im}}}
\def\Coker{\mathop{\mathrm{Coker}}}
\def\Coim{\mathop{\mathrm{Coim}}}

% Boxtimes
%
\DeclareMathSymbol{\boxtimes}{\mathbin}{AMSa}{"02}

%
% Macros for moduli stacks/spaces
%
\def\QCohstack{\mathcal{QC}\!\mathit{oh}}
\def\Cohstack{\mathcal{C}\!\mathit{oh}}
\def\Spacesstack{\mathcal{S}\!\mathit{paces}}
\def\Quotfunctor{\mathrm{Quot}}
\def\Hilbfunctor{\mathrm{Hilb}}
\def\Curvesstack{\mathcal{C}\!\mathit{urves}}
\def\Polarizedstack{\mathcal{P}\!\mathit{olarized}}
\def\Complexesstack{\mathcal{C}\!\mathit{omplexes}}
% \Pic is the operator that assigns to X its picard group, usage \Pic(X)
% \Picardstack_{X/B} denotes the Picard stack of X over B
% \Picardfunctor_{X/B} denotes the Picard functor of X over B
\def\Pic{\mathop{\mathrm{Pic}}\nolimits}
\def\Picardstack{\mathcal{P}\!\mathit{ic}}
\def\Picardfunctor{\mathrm{Pic}}
\def\Deformationcategory{\mathcal{D}\!\mathit{ef}}

% --- appended 2026-10-12 (Chief): the project own definitions of the reference/slogan/history
% display environments, copied mechanically from the shared Stacks preamble lines 8-12
% (\usepackage{verbatim} plus the three \newenvironment lines, which the chapter preamble keeps
% verbatim from the project preamble). Nothing added, removed or redefined: these environments
% carry no mathematical content and the project itself typesets them as comments. Without these
% three lines a counted statement whose own text uses them fails with Environment slogan undefined.
\usepackage{verbatim}
\newenvironment{reference}{\comment}{\endcomment}
\newenvironment{slogan}{\comment}{\endcomment}
\newenvironment{history}{\comment}{\endcomment}
\usepackage{cleveref}
\crefname{originallemma}{original lemma}{original lemmas}
\Crefname{originallemma}{Original lemma}{Original lemmas}
\makeatletter
\newlabel{lemma-quasi-regular-ideal}{{063H}{source}{}{sourcepointer-1}{}}
\newlabel{lemma-quasi-regular-ideal@cref}{{[originallemma][0][]063H}{[1][0][]source}{}{}{}}
\makeatother
\crefname{originallemma}{original lemma}{original lemmas}
\Crefname{originallemma}{Original lemma}{Original lemmas}
\makeatletter
\newlabel{more-algebra-lemma-change-basis}{{0625}{source}{}{sourcepointer-2}{}}
\newlabel{more-algebra-lemma-change-basis@cref}{{[originallemma][0][]0625}{[1][0][]source}{}{}{}}
\makeatother
\title{Elements lifting a basis of the conormal space of a regular ideal generate a regular sequence locally}
\author{The Stacks Project Authors}
\date{stacks-project:divisors.tex:local-snapshot-20261012; GNU FDL 1.2 or later}
\begin{document}
\maketitle
\noindent\textit{Editorial note.} All delivered mathematical text is the author's own, copied from the pinned Stacks Project chapter file pinned-original-source.tex; no statement or proof step was rewritten, reordered or removed.\par\medskip
\noindent\textit{Editorial note.} Following the multi-claim convention for this source, source\_url is the source-level URL https://stacks.math.columbia.edu/ and the source label lemma-generate-regular-ideal is carried in claim\_key and here rather than in a per-claim tag URL.\par\medskip
\noindent\textit{Editorial note.} The preamble is the chapter's shared Stacks Project preamble, retained as sources/stacks-divisors/retained-preamble.tex: a line-for-line excerpt of the shared Stacks preamble staged with this snapshot. No macro was invented or redefined, and this card adds no preamble text of its own.\par\medskip
\noindent\textit{Editorial note.} Statement span 4716-4734 and proof span 4736-4767 come from the halfA pool card operations/pool-cards/stacks-divisors-halfA.json and were re-read in the pinned chapter file before this card was written.\par\medskip
\noindent\textit{Editorial note.} The proof's three parts share one shrinking argument, and the author's remark that the assumptions of the later parts are stronger is preserved.\par\medskip
\noindent\textit{Editorial note.} 2 cross-references occur inside the counted spans; they are listed in external\_references as navigation pointers only.\par\medskip
\noindent\textit{Allowed standard prerequisites (editorial).}\begin{itemize}
\item A quasi-regular sheaf of ideals is locally generated by a quasi-regular sequence, and then its conormal sheaf is free over the quotient with rank the length of that sequence. Cited at proof lines 4736-4742. Produces the reference sequence and shows that the number of given elements equals its length.
\item The Koszul complex of a sequence changes by an invertible change of basis without changing its homology, cited from the more-algebra chapter. Cited at proof lines 4763-4766. Transfers H1-regularity and Koszul-regularity from the reference sequence to the given one.
\item Invertibility of a matrix over a local ring is detected modulo the maximal ideal, so a matrix invertible at a point is invertible on a neighbourhood. Used at proof lines 4747-4752 and again at 4759-4762. Lets the change of basis be inverted after shrinking, and gives that the given elements form a basis of the conormal sheaf locally.
\end{itemize}\medskip
\ifdefined\hypertarget\hypertarget{sourcepointer-1}{}\fi
\noindent\textit{Source reference.} Original lemma 063H: focused-seeds/AGP/sources/stacks-snapshot-20261011/tags.txt, source lines 7831--7831. Same-chapter reference used in the retained author text at source line(s) 4742, 5013. It is supplied only as a navigation pointer so that the retained text compiles; the counted item delivers the author's own statement and the author's complete proof, and no part of the referenced item is copied in or substituted for the counted argument.\par\medskip
\ifdefined\hypertarget\hypertarget{sourcepointer-2}{}\fi
\noindent\textit{Source reference.} Original lemma 0625: focused-seeds/AGP/sources/stacks-snapshot-20261011/tags.txt, source lines 7784--7784. Same-chapter reference used in the retained author text at source line(s) 4764. It is supplied only as a navigation pointer so that the retained text compiles; the counted item delivers the author's own statement and the author's complete proof, and no part of the referenced item is copied in or substituted for the counted argument.\par\medskip
\begin{lemma}
\label{lemma-generate-regular-ideal}
Let $(X, \mathcal{O}_X)$ be a locally ringed space.
Let $\mathcal{J} \subset \mathcal{O}_X$ be a sheaf of ideals.
Let $x \in X$ and $f_1, \ldots, f_r \in \mathcal{J}_x$ whose images
give a basis for the $\kappa(x)$-vector space
$\mathcal{J}_x/\mathfrak m_x\mathcal{J}_x$.
\begin{enumerate}
\item If $\mathcal{J}$ is quasi-regular, then there exists an open
neighbourhood such that $f_1, \ldots, f_r \in \mathcal{O}_X(U)$
form a quasi-regular sequence generating $\mathcal{J}|_U$.
\item If $\mathcal{J}$ is $H_1$-regular, then there exists an open
neighbourhood such that $f_1, \ldots, f_r \in \mathcal{O}_X(U)$
form an $H_1$-regular sequence generating $\mathcal{J}|_U$.
\item If $\mathcal{J}$ is Koszul-regular, then there exists an open
neighbourhood such that $f_1, \ldots, f_r \in \mathcal{O}_X(U)$
form an Koszul-regular sequence generating $\mathcal{J}|_U$.
\end{enumerate}
\end{lemma}

\begin{proof}
First assume that $\mathcal{J}$ is quasi-regular. We may choose an
open neighbourhood $U \subset X$ of $x$ and a quasi-regular sequence
$g_1, \ldots, g_s \in \mathcal{O}_X(U)$ which generates $\mathcal{J}|_U$.
Note that this implies that $\mathcal{J}/\mathcal{J}^2$ is free of
rank $s$ over $\mathcal{O}_U/\mathcal{J}|_U$ (see
Lemma \ref{lemma-quasi-regular-ideal}
and its proof) and hence $r = s$.
We may shrink $U$ and assume $f_1, \ldots, f_r \in \mathcal{J}(U)$.
Thus we may write
$$
f_i = \sum a_{ij} g_j
$$
for some $a_{ij} \in \mathcal{O}_X(U)$. By assumption the matrix
$A = (a_{ij})$ maps to an invertible matrix over $\kappa(x)$.
Hence, after shrinking $U$ once more, we may assume that $(a_{ij})$
is invertible. Thus we see that $f_1, \ldots, f_r$ give a basis
for $(\mathcal{J}/\mathcal{J}^2)|_U$ which proves that $f_1, \ldots, f_r$
is a quasi-regular sequence over $U$.

\medskip\noindent
Note that in order to prove (2) and (3) we may, because the assumptions
of (2) and (3) are stronger than the assumption in (1), already assume that
$f_1, \ldots, f_r \in \mathcal{J}(U)$ and $f_i = \sum a_{ij}g_j$
with $(a_{ij})$ invertible as above, where now $g_1, \ldots, g_r$
is a $H_1$-regular or Koszul-regular sequence. Since the Koszul complex
on $f_1, \ldots, f_r$ is isomorphic to the Koszul complex on
$g_1, \ldots, g_r$ via the matrix $(a_{ij})$ (see
More on Algebra, Lemma \ref{more-algebra-lemma-change-basis})
we conclude that $f_1, \ldots, f_r$ is $H_1$-regular or Koszul-regular
as desired.
\end{proof}

\end{document}
